% !TEX program = lualatex
% ---------------------------------------------------------------------------------------------
%  Readback note 8: Decay of Heisenberg matrix coefficients of homogeneous Fock vectors, uniform in the degree
%
%  Build:  python3 tools/rb.py note <slug>
%
%  One unit of a Lean artifact, rewritten as ordinary mathematics. The prose is hand-written. Lean names,
%  line ranges and counts come from data.tex; quoted Lean statements are copied from the artifact by the build.
%
%  Rules of the text: standard terminology; every symbol is explained where it first appears and listed again
%  in the notation at the end; no letter has two meanings; a proof is complete or says what it takes for granted;
%  Lean names appear only in the LEAN lines and in the concordance.
%
%  The signature rule (the build enforces it): a function, relation or predicate that the artifact defines is
%  written only through its notation macro, which the build makes from the spec (`notation`) with one slot for
%  each explicit argument of the Lean definition. No argument is fixed and dropped, also not "for brevity" and
%  not for one section. Every argument and every hypothesis of a stated declaration stands inside its
%  statement, or in a Setting that the statement cites by number; only implicit and instance arguments (the
%  ambient objects) may stand in the conventions.
% ---------------------------------------------------------------------------------------------
\documentclass[11pt,a4paper]{article}
\ifdefined\pdfvariable\pdfvariable suppressoptionalinfo \numexpr 1+2+4+8\relax\fi
\usepackage[a4paper,hmargin=35mm,top=22mm,bottom=27mm,footskip=12mm,nohead]{geometry}

\usepackage{fontspec}
\directlua{
  luaotfload.add_fallback("gfallback", {"STIXTwoMath-Regular.otf:mode=node;", "DejaVuSans.ttf:mode=node;"})
  luaotfload.add_fallback("gmonofallback", {"DejaVuSansMono.ttf:mode=node;", "STIXTwoMath-Regular.otf:mode=node;"})
}
\setmainfont{NewCM10}[Extension=.otf,UprightFont=*-Book,ItalicFont=*-BookItalic,
  BoldFont=*-Bold,BoldItalicFont=*-BoldItalic,SmallCapsFont=*-Book,SmallCapsFeatures={Letters=SmallCaps},
  RawFeature={fallback=gfallback}]
\setsansfont{NewCMSans10}[Extension=.otf,UprightFont=*-Book,ItalicFont=*-BookOblique,
  BoldFont=*-Bold,BoldItalicFont=*-BoldOblique,RawFeature={fallback=gfallback}]
\setmonofont{JetBrainsMono}[Extension=.otf,UprightFont=*-Regular,ItalicFont=*-Italic,
  BoldFont=*-SemiBold,BoldItalicFont=*-SemiBoldItalic,Scale=MatchLowercase,
  RawFeature={-calt;-liga;fallback=gmonofallback}]
\usepackage{amsmath}
\usepackage{amsthm}
\usepackage[math-style=TeX]{unicode-math}
\setmathfont{NewCMMath-Book.otf}
\emergencystretch=2.5em\relax

\usepackage{xcolor}
\usepackage{tikz}
\usetikzlibrary{calc,arrows.meta,positioning,backgrounds}
\usepackage[protrusion=true,expansion=true]{microtype}
\usepackage{enumitem}
\usepackage{needspace}
\usepackage{booktabs}
\usepackage{array}
\usepackage{longtable}
\usepackage{fancyhdr}
\usepackage{lastpage}
\usepackage{multicol}
\usepackage[hidelinks,pdfusetitle]{hyperref}

\definecolor{ink}{HTML}{0F172A}
\definecolor{inksoft}{HTML}{475569}
\definecolor{inkfaint}{HTML}{94A3B8}
\definecolor{rulegrey}{HTML}{CBD5E1}
\definecolor{arith}{HTML}{C2410C}
\definecolor{est}{HTML}{7C3AED}
\definecolor{ana}{HTML}{0891B2}
\definecolor{ok}{HTML}{047857}
\definecolor{leancol}{HTML}{3730A3}

\title{Decay of Heisenberg matrix coefficients of homogeneous Fock vectors, uniform in the degree}
\author{grh-compact}
\date{}

% ------------------------------------------------------------------------------------- macros
\newcommand{\leanfont}{\ttfamily\fontsize{8.6}{10.5}\selectfont\color{leancol}}
\newcommand{\lean}[1]{\csname LN@#1\endcsname}
\newcommand{\fillin}[1]{{\color{arith}\itshape[#1]}}  % a place of the template that is not written yet; the build reports every one that is left
\newcommand{\leanquote}[1]{\csname LQ@#1\endcsname}
% generated by rb.py note; do not edit
\newcommand{\ConcDecls}{1}
\newcommand{\Pin}{72fead10}
\newcommand{\PinText}{that of commit {\leanfont 72fead10}}
\newcommand{\ArtifactName}{grh-compact}
\expandafter\newcommand\csname LN@normrealInnerrealHeisOple\endcsname{\href{https://github.com/lsernau/grh-lean/blob/72fead102f663d5c7a407964a13dc0cf9d977207/grh-compact/src/New/Weil/FockCoeffDecay.lean\#L13-L35}{{\leanfont norm\_realInner\_realHeisOp\_le}}}
\expandafter\newcommand\csname LQ@normrealInnerrealHeisOple\endcsname{\leanquoteline{0}{lemma norm\_realInner\_realHeisOp\_le \{ι : Type\} [Fintype ι] [Nonempty ι] \{ε P v y\} (σ : USign (n := ι)}\leanquoteline{2}{ε) (hv : v = polyVector ε P) \{d A n\} (hP : ∀ (c : ℝ) x, eval (c • x) P = c \textasciicircum{} d * eval x P) (hA :}\leanquoteline{2}{d + card ι ≤ A) (h0 : A ≠ 0) (hy : 0 < ‖toLp 2 y‖) : ‖realInner v (realHeisOp y (realWeil σ v))‖ ≤}\leanquoteline{2}{‖σ.1.1‖ \textasciicircum{} (card ι + 2 * n) * (realInner v v).re * (4 * (A + n)) \textasciicircum{} n * ‖toLp 2 y‖⁻¹ \textasciicircum{} (2 * n)}}
\newcommand{\ConcordanceRows}{\crow{normrealInnerrealHeisOple}{\href{https://github.com/lsernau/grh-lean/blob/72fead102f663d5c7a407964a13dc0cf9d977207/grh-compact/src/New/Weil/FockCoeffDecay.lean\#L13-L35}{{\leanfont norm\_realInner\_realHeisOp\_le}}}{FockCoeffDecay.lean}{13--35}}
% the notations of the defined functions: one slot for each explicit argument of the Lean declaration
\newcommand{\UnitMembers}{1}
\newcommand{\UnitInNote}{1}
\newcommand{\UnitProofs}{1}
\newcommand{\UnitSteps}{10}
\newcommand{\UnitAllSteps}{10}


\linespread{1.08}                       % 11 pt on 14.7 pt
\setlength{\parindent}{0pt}
\setlength{\parskip}{5.5pt plus 1.5pt minus 0.5pt}
\setlength{\abovedisplayskip}{7.5pt plus 2pt minus 3pt}
\setlength{\belowdisplayskip}{7.5pt plus 2pt minus 3pt}
\setlength{\abovedisplayshortskip}{3pt plus 2pt}
\setlength{\belowdisplayshortskip}{6pt plus 2pt minus 2pt}
\emergencystretch=1.5em
\clubpenalty=10000 \widowpenalty=10000 \displaywidowpenalty=10000
\relpenalty=7000 \binoppenalty=9000
\raggedbottom
\pagestyle{fancy}\fancyhf{}\renewcommand{\headrulewidth}{0pt}
\fancyfoot[C]{\sffamily\fontsize{7.6}{9}\selectfont\color{inkfaint}\thepage\,/\,\pageref*{LastPage}}
\setlist{topsep=2pt,partopsep=0pt,parsep=0pt,itemsep=4pt,leftmargin=2.1em,labelsep=0.55em}

\newcommand{\Q}{\mathbb{Q}}
\newcommand{\Z}{\mathbb{Z}}
\newcommand{\F}{\mathbb{F}}
\newcommand{\len}{\operatorname{len}}
\newcommand{\cok}{\operatorname{cok}}
\newcommand{\End}{\operatorname{End}}
\newcommand{\Hom}{\operatorname{Hom}}
\newcommand{\Gal}{\operatorname{Gal}}
\newcommand{\diag}{\operatorname{diag}}
\newcommand{\kw}[1]{\emph{#1}}
\newcommand{\m}[1]{\mbox{$#1$}}        % a short formula that must not be broken across lines
\newcommand{\head}[1]{\par\addvspace{17pt plus 4pt minus 2pt}\Needspace{6\baselineskip}%
  \noindent{\fontsize{13.2}{16}\selectfont\bfseries\color{ink}#1}\par\nopagebreak\vspace{-1pt}}
\newcommand{\sect}[2]{\refstepcounter{section}\label{#2}\par\addvspace{19pt plus 4pt minus 2pt}\Needspace{7\baselineskip}%
  \noindent\llap{{\fontsize{13.2}{16}\selectfont\bfseries\color{inkfaint}\thesection}\hspace{9pt}}%
  {\fontsize{13.2}{16}\selectfont\bfseries\color{ink}#1}\par\nopagebreak\vspace{-1pt}}
\newcommand{\subhead}[1]{\par\addvspace{9pt plus 2pt}\Needspace{4\baselineskip}\noindent\textbf{#1}\enspace\ignorespaces}
% where a statement lives in the Lean archive
\newcommand{\leantag}{{\sffamily\bfseries\fontsize{6.3}{8}\selectfont\color{leancol}\addfontfeatures{LetterSpace=7}LEAN}}
\newcommand{\sep}{\unskip\nobreak\hspace{5pt}{\color{inkfaint}\textperiodcentered}\hspace{5pt plus 2pt}\ignorespaces}
\newcommand{\leanline}[1]{\par\nopagebreak\vspace{-2.5pt}{\fontsize{8.9}{12.4}\selectfont\color{inksoft}\raggedright
  \noindent\llap{\leantag\hspace{9pt}}#1\par}\vspace{1pt}}
% a displayed statement in a tinted box: #1 colour, #2 text
\newcommand{\statement}[2]{\par\addvspace{7pt}\begingroup\setlength{\fboxsep}{8.5pt}\setlength{\fboxrule}{0.6pt}%
  \noindent\fcolorbox{#1!30!white}{#1!3!white}{\parbox{\dimexpr\linewidth-2\fboxsep-2\fboxrule\relax}{#2}}%
  \endgroup\par\addvspace{7pt}}
\newcommand{\nt}[2]{\par\noindent\hangindent=1.1em\hangafter=1 #1\enspace{\color{inksoft}#2}\par\vspace{1pt}}
\newcommand{\legdot}[1]{\tikz[baseline=-0.62ex]\fill[#1] (0,0) circle (2.5pt);\hspace{3pt}}
\newcommand{\crow}[4]{\csname CS@#1\endcsname & #2 & \textcolor{inksoft}{#3} & \textcolor{inksoft}{#4}\\}
\newcommand{\leanquoteline}[2]{\par\noindent\hangindent=\dimexpr#1\fontdimen2\font+2.2em\relax\hangafter=1
  \hspace*{#1\fontdimen2\font}#2}
\newcommand{\quoted}[1]{\par\vspace{3pt}{\ttfamily\fontsize{7.5}{10.4}\selectfont\color{leancol}\raggedright
  \setlength{\parskip}{0pt}\leanquote{#1}\par}\vspace{4pt}}
\newcommand{\bx}[3]{{\sffamily\bfseries\fontsize{7.3}{9.2}\selectfont\color{#1!82!black}#2}\\[1.5pt]#3}
\newcommand{\cs}[2]{\expandafter\newcommand\csname CS@#1\endcsname{#2}}

% theorems
\newtheoremstyle{plainx}{7pt plus 2pt}{7pt plus 2pt}{\itshape}{}{\bfseries}{.}{0.5em}{}
\newtheoremstyle{uprightx}{7pt plus 2pt}{7pt plus 2pt}{}{}{\bfseries}{.}{0.5em}{}
\theoremstyle{plainx}
\newtheorem{theorem}{Theorem}[section]
\newtheorem{lemma}[theorem]{Lemma}
\newtheorem{proposition}[theorem]{Proposition}
\newtheorem{corollary}[theorem]{Corollary}
\theoremstyle{uprightx}
\newtheorem{definition}[theorem]{Definition}
\newtheorem{setting}[theorem]{Setting}          % hypotheses shared by several statements; each of them cites it by number
\newtheorem{remark}[theorem]{Remark}
\newtheorem{example}[theorem]{Example}
\renewcommand{\qedsymbol}{{\setlength{\fboxsep}{0pt}\setlength{\fboxrule}{0.5pt}\fbox{\rule{0pt}{0.46em}\hspace{0.46em}}}}
% what holds throughout the note: the ambient objects only (implicit and instance arguments), never a hypothesis
\newenvironment{conventions}{\par\addvspace{5pt}\noindent\textbf{Conventions.}\enspace\ignorespaces}{\par\addvspace{3pt}}
\makeatletter
\renewenvironment{proof}[1][Proof]{\par\pushQED{\qed}\normalfont\topsep0pt\partopsep0pt
  \trivlist\item[\hskip\labelsep\itshape #1.]\ignorespaces}{\popQED\endtrivlist\@endpefalse}
\makeatother

% where each declaration of the concordance is treated: \cs{key}{place}, the place naming a \label or a lettered item
\cs{normrealInnerrealHeisOple}{Theorem 1}
% notation of this note (objects of the background library are written with all their arguments)
\newcommand{\CC}{\mathbb{C}}
\newcommand{\RR}{\mathbb{R}}
\newcommand{\Fock}[1]{\mathcal{F}(#1)}                 % the Fock space of signature #1
\newcommand{\heis}[2]{\varrho_{#1}(#2)}                % the Heisenberg operator of the vector #2
\newcommand{\weil}[2]{\omega_{#1}(#2)}                 % the oscillator representation at #2
\newcommand{\inn}[3]{\langle #2,#3\rangle_{#1}}        % the inner product of the Fock space
\newcommand{\pv}[2]{v_{#1,#2}}                         % the Fock vector of the polynomial #2
\newcommand{\hol}[2]{\operatorname{hol}_{#1}(#2)}      % conjugation of the positive coordinates
\newcommand{\vac}{\varphi_0}                           % the Gaussian
\newcommand{\bich}[3]{\beta_{#1}(#2,#3)}               % the bicharacter
\newcommand{\mult}[2]{\Delta_{#1}(#2)}                 % the multiplier
\newcommand{\opn}[1]{\lVert #1\rVert}                  % operator norm
\newcommand{\lone}[1]{\lVert #1\rVert_{1}}             % L^1 norm
\newcommand{\ee}{\mathrm{e}}
\newcommand{\dd}{\mathrm{d}}
\newcommand{\dg}{\mathrm{dg}}                           % the degree of the homogeneous polynomial
\newcommand{\lnm}[1]{{\leanfont #1}}

\begin{document}
\color{ink}

{\raggedright
{\sffamily\bfseries\fontsize{7.2}{9}\selectfont\color{leancol}\addfontfeatures{LetterSpace=8}READBACK NOTE 8\par}
\vspace{5pt}
{\fontsize{21.5}{25.5}\selectfont\bfseries Decay of Heisenberg matrix coefficients of homogeneous Fock vectors, uniform in the degree\par}
\vspace{5pt}
{\fontsize{12.4}{16}\selectfont\itshape\color{inksoft} One unit of the Lean artifact \ArtifactName, read back into
mathematics\par}}
\vspace{8pt}
{\color{rulegrey}\hrule height 0.6pt}
\vspace{7pt}
\thispagestyle{empty}

\textbf{Summary.} In the Fock model of the oscillator representation of a unitary group $U(\varepsilon)$ of any
signature, let $v$ be the vector of a homogeneous polynomial. The matrix coefficient of~$v$ against its image
under a group element followed by a Heisenberg translation by a vector~$y$ decays faster than every power
of~$|y|$. The unit proves this with explicit constants that depend on the degree of the polynomial only through
a bound \m{A\ge\text{degree}+\text{number of variables}}: the decay of order \m{|y|^{-2\nu}} costs the factor
\m{(4(A+\nu))^{\nu}}. This is the sense of ``uniform in the degree'' in the title: no feature of the polynomial
other than~$A$ enters the constants. The bound itself is not independent of the degree: for fixed~$\nu$ the factor
grows like the $\nu$-th power of~$A$. The proof writes the coefficient as an integral, compares $v$ with a dilate of itself, and
uses three elementary inequalities. This is the first of three notes on the theta correspondence behind
Theorem~C; it is used where a theta kernel is compared with a matrix coefficient.

\statement{leancol}{\fontsize{9.6}{13.2}\selectfont\textbf{What this note is.} The artifact \ArtifactName{} is a
body of Lean code: the compact edition of a formal proof, due to L.~Alpöge with the AI system Claude, whose final
theorems are the Riemann hypothesis and its generalisation to Dirichlet \emph{L}-functions. That headline is a claim.
The authors report that every file compiles and that Lean's kernel accepted the proofs; an independent build is
recorded in the handout's record of checks when it reports. No mathematician has read the argument through, and
nothing here says that the Riemann hypothesis is proved. This note treats one unit of the artifact,
which consists of a single lemma (\UnitInNote{} of \UnitMembers{} declarations). Its proof is given in full. The Fock model
itself (the space, the Heisenberg operators, the oscillator representation) belongs to the background library of
the artifact and is quoted in §\ref{s:setting}, with the few properties that the proof uses. The statement was
compared with its Lean text, argument by argument. The Lean text read is \PinText; Appendix~A is the concordance
of \ConcDecls{} declaration.}

\sect{The statement}{s:intro}
The objects are defined in §\ref{s:setting}: $U(\varepsilon)$ is the unitary group of the diagonal hermitian form
with entries~$\varepsilon$, the space $\Fock{\varepsilon}$ is the Fock space of signature~$\varepsilon$ with its
inner product \m{\inn{\varepsilon}{\,\cdot\,}{\,\cdot\,}}, the operators \m{\heis{\varepsilon}{y}} and
\m{\weil{\varepsilon}{u}} are the Heisenberg translation by~$y$ and the oscillator representation at~$u$, and
\m{\pv{\varepsilon}{p}} is the vector of the polynomial~$p$. Norms of vectors of $\CC^I$ are Euclidean, and
\m{\opn{u}} is the operator norm of the matrix~$u$ for these norms.

\statement{arith}{\phantomsection\label{t:1}\textbf{Theorem 1} (decay of matrix coefficients, with explicit
dependence on the degree).\enspace \emph{Let $I$ be a non-empty finite set, \m{\varepsilon\colon I\to\{1,-1\}} a family of signs,
\m{u\in U(\varepsilon)}, $p$ a complex polynomial in the variables $X_i$, \m{i\in I}, and $v$ an element of the
Fock space $\Fock{\varepsilon}$ with \m{v=\pv{\varepsilon}{p}}. Let $\dg$, $A$ and $\nu$ be natural numbers and
let \m{y\in\CC^I}. Suppose that $p$ is homogeneous of degree~$\dg$, in the sense that
\m{p(cz)=c^{\dg}p(z)} for every real number~$c$ and every \m{z\in\CC^I}; that \m{\dg+|I|\le A}; that
\m{A\ne0}; and that \m{|y|>0}. Then}
\[
\bigl|\inn{\varepsilon}{v}{\heis{\varepsilon}{y}\,\weil{\varepsilon}{u}\,v}\bigr|
\ \le\ \opn{u}^{\,|I|+2\nu}\cdot\operatorname{Re}\inn{\varepsilon}{v}{v}\cdot\bigl(4(A+\nu)\bigr)^{\nu}\cdot\bigl(|y|^{-1}\bigr)^{2\nu}.
\]}
\leanline{\lean{normrealInnerrealHeisOple}}

The number \m{\inn{\varepsilon}{v}{v}} is the integral of \m{|v|^2}, so it is real and its real part is itself.
For \m{\nu=0} the bound reads \m{\opn{u}^{|I|}\inn{\varepsilon}{v}{v}}, which is weaker than the bound
\m{\inn{\varepsilon}{v}{v}} that the Cauchy--Schwarz inequality gives together with the unitarity in (Q6), a fact
this note does not use; the content lies in large~$\nu$. If \m{\nu\le|y|^2/(4\ee\opn{u}^2)-A}, then
\m{4(A+\nu)\opn{u}^2|y|^{-2}\le\ee^{-1}} and the right side is at most
\m{\opn{u}^{|I|}\inn{\varepsilon}{v}{v}\,\ee^{-\nu}}; so the theorem gives a bound of Gaussian type in~$|y|$. What
matters in the application is that the constant \m{(4(A+\nu))^{\nu}} depends on the polynomial only through the
bound~$A$ for its degree plus the number of variables. The theorem is a quantitative form, with the dependence on
the degree made explicit, of the rapid decay in~$y$ of such coefficients.

\sect{The setting}{s:setting}
\begin{conventions}\phantomsection\label{c:conventions}
The natural numbers are \m{0,1,2,\dots}. Throughout, $I$ is a non-empty finite set and \m{\varepsilon\colon
I\to\{1,-1\}} a family of signs; \m{|I|} is the number of elements of~$I$. Vectors of $\CC^I$ are columns indexed
by~$I$, with the Euclidean norm \m{|z|=(\sum_{i\in I}|z_i|^2)^{1/2}}, and for \m{z,w\in\CC^I} we put
\m{w^{*}\varepsilon z=\sum_{i\in I}\overline{w_i}\,\varepsilon_i\,z_i}. Integrals over $\CC^I$ are taken with
respect to the Lebesgue measure divided by \m{\pi^{|I|}}, written~\m{\dd z}; so
\m{\int_{\CC^I}\exp(-|z|^2)\,\dd z=1}, and the substitution \m{z\mapsto z/c}, for a real number \m{c>0}, multiplies
integrals by \m{c^{2|I|}}. For a
square matrix~$u$ indexed by~$I$, \m{\opn{u}} is its operator norm on $\CC^I$ with the Euclidean norm. The
objects of this section are those of the background library of the artifact; they are recalled here because the
statement of Theorem~1 uses them. Letters of the handout's overview that are not displayed with the data as an
argument (among them $\varepsilon$, $\beta$, $\nu$, $w$ and~$j$) have, in this note, only the meaning given in its
notation list; the degree is written~$\dg$, because $\delta$ is the bound of Theorem~C in the overview.
\end{conventions}

\begin{definition}[the Heisenberg bicharacter and the Fock space]\label{d:fock}
For \m{z,w\in\CC^I} put
\[
\bich{\varepsilon}{z}{w}=\exp\Bigl(\tfrac12\bigl(w^{*}\varepsilon z-z^{*}\varepsilon w\bigr)\Bigr)
=\exp\bigl(\mathrm{i}\operatorname{Im}(w^{*}\varepsilon z)\bigr),\qquad \vac(z)=\exp\bigl(-\tfrac12|z|^2\bigr).
\]
So \m{\beta_{\varepsilon}} has absolute value~$1$, is multiplicative in each variable, and
\[\bich{\varepsilon}{w}{z}=\bich{\varepsilon}{z}{w}^{-1}=\overline{\bich{\varepsilon}{z}{w}};\] the Gaussian
$\vac$ is real, even, and \m{0<\vac\le1}. The \kw{Fock space} $\Fock{\varepsilon}$ is the space of all integrable
functions \m{v\colon\CC^I\to\CC} that satisfy the reproducing identity
\[
v(w)=\int_{\CC^I}v(z)\,\bich{\varepsilon}{z}{w}\,\vac(w-z)\,\dd z\qquad\text{for every } w\in\CC^I,
\]
with the norm \m{\lone{v}=\int|v(z)|\,\dd z}. For \m{v,\mathbf{w}\in\Fock{\varepsilon}} the \kw{inner product} is
\[
\inn{\varepsilon}{v}{\mathbf{w}}=\int_{\CC^I}\overline{v(z)}\,\mathbf{w}(z)\,\dd z ;
\]
it is conjugate-linear in the first variable. For
\m{y\in\CC^I} the \kw{Heisenberg operator} \m{\heis{\varepsilon}{y}} sends \m{v\in\Fock{\varepsilon}} to the function
\[
\bigl(\heis{\varepsilon}{y}v\bigr)(z)=\bich{\varepsilon}{y}{z}\,v(z-y).
\]
\end{definition}
\leanline{in the background library: \lnm{realHeis}, \lnm{realVac}, \lnm{realMeasure}, \lnm{fockSpace},
\lnm{realInner}, \lnm{realHeisOp}\sep the general construction is \lnm{Heisenberg.IsVacuum.Model} with
\lnm{heisOp} and \lnm{modelInner}}

The elements of $\Fock{\varepsilon}$ are functions, not classes of functions: the reproducing identity determines
the values. It also gives \m{|v(w)|\le\lone{v}} for all~$w$, because \m{|\beta_{\varepsilon}|=1} and
\m{\vac\le1}. So every \m{v\in\Fock{\varepsilon}} is bounded, \m{|v|^2\le\lone{v}\,|v|} is integrable, and the inner
product is defined and bounded: \m{|\inn{\varepsilon}{v}{\mathbf{w}}|\le\lone{v}\lone{\mathbf{w}}}. The Heisenberg operators map
$\Fock{\varepsilon}$ to itself, preserve the norm and the inner product, and satisfy
\begin{equation}\label{e:heis}
\heis{\varepsilon}{y}\,\heis{\varepsilon}{w}=\bich{\varepsilon}{y}{w}\,\heis{\varepsilon}{y+w}\qquad(y,w\in\CC^I).
\end{equation}
The Gaussian $\vac$ belongs to $\Fock{\varepsilon}$; it is called the \kw{vacuum vector}. For \m{\varepsilon=1}
the functions \m{z\mapsto v(z)/\vac(z)}, \m{v\in\Fock{\varepsilon}}, are antiholomorphic, and one has the
integrable functions of the classical Fock space; in general the positive coordinates enter through their
conjugates and the negative ones directly, as the next definition shows. The background library works with this
space of integrable functions and its norm of integrable functions, and not with the Hilbert space completion.

\begin{definition}[polynomial vectors]\label{d:poly}
For \m{z\in\CC^I} let \m{\hol{\varepsilon}{z}\in\CC^I} be the vector with coordinates \m{\overline{z_i}} for
\m{\varepsilon_i=1} and \m{z_i} for \m{\varepsilon_i=-1}. For a complex polynomial~$p$ in the variables $X_i$,
\m{i\in I}, the \kw{Fock vector of~$p$} is the function
\[
\pv{\varepsilon}{p}(z)=p\bigl(\hol{\varepsilon}{z}\bigr)\,\vac(z)\qquad(z\in\CC^I).
\]
\end{definition}
\leanline{in the background library: \lnm{holCoord}, \lnm{polyVector}}

\begin{definition}[the unitary group, its multiplier, the oscillator representation]\label{d:weil}
The group $U(\varepsilon)$ consists of the invertible complex matrices~$u$ indexed by~$I$ with
\m{u^{*}\operatorname{diag}(\varepsilon)\,u=\operatorname{diag}(\varepsilon)}, where $u^{*}$ is the conjugate
transpose; it acts on $\CC^I$ by \m{z\mapsto uz}. For \m{u\in U(\varepsilon)} let $u_{--}$ be the square block
of~$u$ whose rows and columns are the indices~$i$ with \m{\varepsilon_i=-1}, and put
\m{\mult{\varepsilon}{u}=\overline{\det u_{--}}}, the \kw{multiplier} (it is~$1$ if no sign is~$-1$). The
\kw{oscillator representation} of \m{u\in U(\varepsilon)} is the operator on $\Fock{\varepsilon}$
\[
\weil{\varepsilon}{u}\,v=\mult{\varepsilon}{u}\int_{\CC^I}v(z)\;\heis{\varepsilon}{uz}\,\vac\;\dd z ,
\]
an integral of a function with values in the Banach space $\Fock{\varepsilon}$.
\end{definition}
\leanline{in the background library: \lnm{USign}, \lnm{realMultiplier}, \lnm{realAction}, \lnm{realWeil}\sep the
general construction is \lnm{Heisenberg.IsVacuum.weilOp}, \lnm{weilRep}}

\statement{inksoft}{\textbf{Quoted from the background library.} These facts about the Fock model are taken for
granted. Facts (Q2) to (Q5) have proofs of a line, which are indicated; (Q1) and (Q6) are part of the classical
theory of the Fock model of the oscillator representation of unitary groups of any signature (references in §\ref{s:notes}, note~5).

\vspace{3pt}\textbf{(Q1)} For every complex polynomial~$p$ in the variables $X_i$, \m{i\in I}, the function
\m{\pv{\varepsilon}{p}} belongs to $\Fock{\varepsilon}$; in particular it is integrable and satisfies the reproducing
identity.

\vspace{3pt}\textbf{(Q2)} \emph{Reproducing property.} For \m{v\in\Fock{\varepsilon}} and \m{w\in\CC^I}:
\m{\inn{\varepsilon}{\heis{\varepsilon}{w}\vac}{v}=v(w)}. Indeed the left side is
\m{\int\overline{\bich{\varepsilon}{w}{z}}\,\vac(z-w)\,v(z)\,\dd z}, which is the right side of the reproducing
identity because \m{\overline{\bich{\varepsilon}{w}{z}}=\bich{\varepsilon}{z}{w}} and $\vac$ is real and even.

\vspace{3pt}\textbf{(Q3)} For \m{u\in U(\varepsilon)}: \m{\bich{\varepsilon}{uz}{uw}=\bich{\varepsilon}{z}{w}} for all
\m{z,w\in\CC^I}, and \m{z\mapsto uz} preserves the measure~\m{\dd z}. The first holds because
\m{(uw)^{*}\varepsilon(uz)=w^{*}\varepsilon z}, the second because \m{|\det u|=1}.

\vspace{3pt}\textbf{(Q4)} For \m{u\in U(\varepsilon)}: \m{\opn{u}\ge1} and
\m{|\mult{\varepsilon}{u}|\le\opn{u}^{|I|}}. Indeed \m{u^{-1}=\operatorname{diag}(\varepsilon)\,u^{*}\operatorname{diag}(\varepsilon)}
has the norm of~$u$, so \m{1=\opn{u^{-1}u}\le\opn{u}^2}; and a principal minor of size~$j$ of a matrix~$u$ has absolute
value at most \m{\opn{u}^{j}\le\opn{u}^{|I|}}.

\vspace{3pt}\textbf{(Q5)} A bounded linear map~$T$ from $\Fock{\varepsilon}$ to~$\CC$ commutes with the integral of
Definition~\ref{d:weil}: \m{T\bigl(\int v(z)\,\heis{\varepsilon}{uz}\vac\,\dd z\bigr)=\int v(z)\,T(\heis{\varepsilon}{uz}\vac)\,\dd z}.

\vspace{3pt}\textbf{(Q6)} \emph{Not used in this note.} The map \m{u\mapsto\weil{\varepsilon}{u}} is a
homomorphism from $U(\varepsilon)$ to the bounded invertible operators of $\Fock{\varepsilon}$, it preserves the
inner product, and \m{\weil{\varepsilon}{u}\,\heis{\varepsilon}{y}=\heis{\varepsilon}{uy}\,\weil{\varepsilon}{u}}.
This is the statement that the oscillator representation splits over $U(\varepsilon)$ with the multiplier
\m{\Delta_{\varepsilon}}.}
\leanline{in the background library: (Q1) is part of \lnm{polyVector}\sep (Q2) \lnm{Heisenberg.IsVacuum.modelInner\_heisOp\_vacElem}\sep
(Q3) is part of \lnm{realAction}\sep (Q4) \lnm{U.one\_le\_norm}, \lnm{norm\_realMultiplier\_le}\sep
(Q5) \lnm{Heisenberg.Rep.map\_integrated\_eq\_integral}\sep (Q6) is part of \lnm{realAction}, with \lnm{weilRep},
\lnm{modelInner\_weilRep}, \lnm{weilRep\_heisOp}}

\sect{Proof of Theorem 1}{s:proofs}
Three lemmas prepare the proof. In all of them $I$ and $\varepsilon$ are as in the conventions.

\begin{lemma}[the coefficient as an integral]\label{l:integral}
Let \m{u\in U(\varepsilon)}, \m{v\in\Fock{\varepsilon}} and \m{y\in\CC^I}. Then
\begin{align*}
\inn{\varepsilon}{v}{\heis{\varepsilon}{y}\,\weil{\varepsilon}{u}\,v}
&=\mult{\varepsilon}{u}\int_{\CC^I}v(z)\,\bich{\varepsilon}{y}{uz}\,\overline{v(uz+y)}\;\dd z ,\qquad\text{and so}\\
\bigl|\inn{\varepsilon}{v}{\heis{\varepsilon}{y}\,\weil{\varepsilon}{u}\,v}\bigr|
&\le\opn{u}^{|I|}\int_{\CC^I}|v(z)|\,|v(uz+y)|\;\dd z .
\end{align*}
\end{lemma}
\begin{proof}
The map \m{T\colon \mathbf{w}\mapsto\inn{\varepsilon}{v}{\heis{\varepsilon}{y}\mathbf{w}}} is linear and bounded on
$\Fock{\varepsilon}$. By Definition~\ref{d:weil} and (Q5),
\[
T\bigl(\weil{\varepsilon}{u}v\bigr)=\mult{\varepsilon}{u}\int v(z)\;\inn{\varepsilon}{v}{\heis{\varepsilon}{y}\,\heis{\varepsilon}{uz}\,\vac}\;\dd z .
\]
By \eqref{e:heis} and by (Q2), for every \m{w\in\CC^I},
\[
\heis{\varepsilon}{y}\,\heis{\varepsilon}{uz}\,\vac=\bich{\varepsilon}{y}{uz}\,\heis{\varepsilon}{uz+y}\,\vac ,\qquad
\inn{\varepsilon}{v}{\heis{\varepsilon}{w}\vac}=\overline{\inn{\varepsilon}{\heis{\varepsilon}{w}\vac}{v}}=\overline{v(w)} .
\]
This is the identity. The inequality follows from \m{|\beta_{\varepsilon}|=1} and (Q4).
\end{proof}

\begin{lemma}[dilation of the vector of a homogeneous polynomial]\label{l:dilate}
Let $p$ be a complex polynomial in the variables $X_i$, \m{i\in I}, and $\dg$ a natural number with
\m{p(cz)=c^{\dg}p(z)} for every real~$c$ and every \m{z\in\CC^I}. Then for every real \m{c>0} and every
\m{z\in\CC^I}
\[
\bigl|\pv{\varepsilon}{p}(z)\bigr|=c^{\dg}\,\exp\Bigl(-\tfrac12\bigl(1-c^{-2}\bigr)|z|^2\Bigr)\,\bigl|\pv{\varepsilon}{p}(z/c)\bigr| .
\]
\end{lemma}
\begin{proof}
The map \m{z\mapsto\hol{\varepsilon}{z}} is linear over the real numbers, so
\[
p\bigl(\hol{\varepsilon}{z/c}\bigr)=c^{-\dg}\,p\bigl(\hol{\varepsilon}{z}\bigr),\qquad
\vac(z/c)=\vac(z)\exp\bigl(\tfrac12(1-c^{-2})|z|^2\bigr).
\]
Multiply the two.
\end{proof}

\begin{lemma}[three elementary inequalities]\label{l:elementary}
\textup{(i)} \m{1+B\le\ee^{B}} for real~$B$. \textup{(ii)} \m{\ee^{1-B}\le B^{-1}} for real \m{B>0}.
\textup{(iii)} For \m{u\in U(\varepsilon)} and \m{y,z\in\CC^I}:
\[
|y|^2\le2\opn{u}^2\bigl(|z|^2+|uz+y|^2\bigr).
\]
\end{lemma}
\begin{proof}
(i) is the convexity of the exponential function, and (ii) is (i) applied to \m{B-1}: \m{B\le\ee^{B-1}}. For
(iii), the triangle inequality and (Q4) give
\[
|y|\le|uz+y|+\opn{u}|z|\le\opn{u}\bigl(|uz+y|+|z|\bigr),
\]
and \m{(\alpha+\gamma)^2\le2(\alpha^2+\gamma^2)} for real \m{\alpha,\gamma}.
\end{proof}

\begin{proof}[Proof of Theorem 1]
Since $I$ is not empty and \m{\dg+|I|\le A}, we have \m{A\ge1}. Put \m{c=\sqrt{\nu/A+1}}, so that \m{c\ge1} and
\m{1-c^{-2}=\nu/(\nu+A)}. By (Q1), \m{v=\pv{\varepsilon}{p}} lies in $\Fock{\varepsilon}$, so $|v|^2$ is
integrable; hence so are \m{z\mapsto|v(z/c)|^2} and, by (Q3), \m{z\mapsto|v((uz+y)/c)|^2}, with
\begin{equation}\label{e:two}
\int_{\CC^I}|v(z/c)|^2\,\dd z=\int_{\CC^I}\bigl|v\bigl((uz+y)/c\bigr)\bigr|^2\,\dd z=c^{2|I|}\,\inn{\varepsilon}{v}{v}.
\end{equation}
Fix \m{z\in\CC^I} and put \m{B=(|z|^2+|uz+y|^2)/(2(\nu+A))}. If $B$ were~$0$, then \m{z=0} and \m{uz+y=0}, so
\m{y=0}, which is excluded; hence \m{B>0}. Lemma~\ref{l:dilate}, applied at the points $z$ and \m{uz+y}, gives
\[
|v(z)|\,|v(uz+y)|=c^{2\dg}\,\ee^{-\nu B}\;|v(z/c)|\,\bigl|v\bigl((uz+y)/c\bigr)\bigr| .
\]
Now \m{c^2=1+\nu/A\le\ee^{\nu/A}} by Lemma~\ref{l:elementary}(i), and \m{(\dg+|I|)/A\le1}, so
\[c^{2\dg}=c^{2(\dg+|I|)}\,c^{-2|I|}\le\ee^{\nu}\,c^{-2|I|}.\]
By Lemma~\ref{l:elementary}(ii) and~(iii),
\[
\ee^{\nu}\,\ee^{-\nu B}=\bigl(\ee^{1-B}\bigr)^{\nu}\le B^{-\nu}
\le\Bigl(\frac{4(A+\nu)\opn{u}^2}{|y|^2}\Bigr)^{\nu}.
\]
With \m{2\alpha\gamma\le\alpha^2+\gamma^2} for the last two factors this gives, for every~$z$,
\[
|v(z)|\,|v(uz+y)|\ \le\ \frac{1}{2\,c^{2|I|}}\Bigl(\frac{4(A+\nu)\opn{u}^2}{|y|^2}\Bigr)^{\nu}
\Bigl(|v(z/c)|^2+\bigl|v\bigl((uz+y)/c\bigr)\bigr|^2\Bigr).
\]
Integrate over~$z$ and use \eqref{e:two}: the integral of the left side is at most
\[
\bigl(4(A+\nu)\bigr)^{\nu}\,\opn{u}^{2\nu}\,|y|^{-2\nu}\,\inn{\varepsilon}{v}{v}.
\]
Lemma~\ref{l:integral} multiplies this by \m{\opn{u}^{|I|}}.
\end{proof}
\leanline{\lean{normrealInnerrealHeisOple}\sep the formal proof is one chain of the same steps: the constant~$c$,
the two integrable functions, the integral formula, the bound of the multiplier, the pointwise inequality}

\sect{Where the result is used}{s:use}
The theorem has one consumer, in the unit on the rows (the forms on the indefinite unitary group). There a theta
kernel for a pair of unitary groups, formed with a datum that is a character on the matrices modulo a
modulus, is compared with the matrix coefficient \m{\inn{\varepsilon}{v}{\weil{\varepsilon}{u}v}} of a
homogeneous polynomial vector~$v$ for real points~$u$ of the two groups. The difference is a sum, over the non-zero
vectors~$y$ of a lattice, of the coefficients \m{\inn{\varepsilon}{v}{\heis{\varepsilon}{y}\weil{\varepsilon}{u}v}}.
Theorem~1 is applied to each term, with \m{A=\dg+|I|}; the vectors~$y$ have norm at least a quantity fixed by
the first minimum of the lattice, and the sum of \m{|y|^{-2\nu}} over the lattice converges when $2\nu$ exceeds
its real rank. The outcome is a bound for the difference whose dependence on the degree~$\dg$ is explicit, through the factor
\m{(4(\dg+|I|+\nu))^{\nu}}, polynomial in~$\dg$ for fixed~$\nu$; this is what allows degrees that grow with the
data. That unit states the result as an Input.

\sect{Notes on the formal text}{s:notes}
\begin{enumerate}[label=\textup{\arabic*.},leftmargin=1.6em,itemsep=5pt]
\item \emph{A hypothesis that follows from the others.} The formal statement assumes \m{A\ne0}. Since $I$ is not
empty and \m{\dg+|I|\le A}, this follows from the other hypotheses; the formal proof uses it where the
proof above uses \m{A\ge1}, namely for the identity \m{1-c^{-2}=\nu/(\nu+A)}. Theorem~1 keeps the hypothesis because
the formal statement has it.
\item \emph{Homogeneity.} The formal statement does not use the notion of a homogeneous polynomial of its
library. It asks for \m{p(cz)=c^{\dg}p(z)} for all real~$c$, zero and negative values included, and all
\m{z\in\CC^I}. For a polynomial this is the same as homogeneity of degree~$\dg$; only positive~$c$ is used.
\item \emph{The form of the bound.} The right side is written in the formal text as the product
\m{\opn{u}^{|I|+2\nu}\cdot\operatorname{Re}\inn{\varepsilon}{v}{v}\cdot(4(A+\nu))^{\nu}\cdot(|y|^{-1})^{2\nu}},
with the real part and with the inverse of~$|y|$ raised to the power~$2\nu$; Theorem~1 reproduces this form. The
vector~$v$ is an argument of its own, tied to~$p$ by the hypothesis \m{v=\pv{\varepsilon}{p}}.
\item \emph{Steps carried out by automation.} The reader of the unit marked the last four steps of the
structured proof (the dilation identity, the collection of powers, the two chains of inequalities) as claims
about the mathematical content, not replayed goal by goal. I read these lines in the Lean text. The dilation
identity is stated there explicitly, as
\m{|v(z/c)|=|v(z)|\exp(|z|^2/(\nu+A)/2)^{\nu}/|c|^{\dg}}, which is Lemma~\ref{l:dilate} for the chosen~$c$. The
chains then use, in this order: the square of the square root, \m{1+B\le\ee^{B}}, the hypothesis
\m{\dg+|I|\le A}, \m{A/A\le1}, \m{\ee\cdot B\le\ee^{B}}; then the triangle inequality, the bound of \m{|uz|} by
\m{\opn{u}|z|}, \m{\opn{u}\ge1}, and \m{(\alpha+\gamma)^2\le2(\alpha^2+\gamma^2)}. These are the steps of the proof
above. What the automation proves between them was not replayed by me either: no Lean was run for this note.
\item \emph{A reading of the dossier that was corrected after this note was written.} The first version of the
structured proof described the second integrable function as \m{z\mapsto|v(uz+y)|^2} and said that it integrates
to \m{\inn{\varepsilon}{v}{v}}. In the Lean text it is the composite of \m{z\mapsto|v(z/c)|^2} with the
measure-preserving map \m{z\mapsto uz+y}, that is \m{z\mapsto|v((uz+y)/c)|^2}, and both functions integrate to
\m{c^{2|I|}\inn{\varepsilon}{v}{v}}, as in \eqref{e:two}. With the dilated second function the constants match
exactly; with the undilated one the two integrals would have different scales. The second reader of the artifact,
who read the unit independently, reported the same reading of these two lines, and the structured proof now
follows it in the steps that name the two functions, the pointwise bound and the two integrals. The quotations
of the dossier were exact throughout; only its words were affected.
\item \emph{References for the quoted facts.} The catalogue of the handout records no confirmed place for the
facts (Q1) to (Q6): \emph{reference to be confirmed}. The files of the background library name their sources in
their headers, for the Fock model of the oscillator representation of the indefinite unitary groups: M.~G.~Davidson, \emph{The
harmonic representation of U(p,q) and its connection with the generalized unit disk}, Pacific J.~Math.~129
(1987), §1 and Prop.~2.2; G.~B.~Folland, \emph{Harmonic Analysis in Phase Space} (1989), ch.~1 and~4; H.~Reiter,
\emph{Metaplectic Groups and Segal Algebras}, Lecture Notes in Math.~1382 (1989). The artifact's own list of
references marks the first as ``work found, no place given'' and the second as ``place confirmed at second
hand''. The same headers say how the formal text differs from these sources: the space is the space of
integrable functions with the reproducing identity, not the Hilbert space, and no kernel formula for
\m{\weil{\varepsilon}{u}} is proved. For (Q6) the catalogue notes that the formal definition of the action
contains the proof that the cocycle is trivial, by a Gaussian integral. The second reader located the classical
statement in Davidson's paper, §1, formula (1.9a), Lemma~1.11 and Theorem~1.12, pp.~39--40, where the phases of
the Fock intertwiners are fixed and shown to give a continuous representation of the indefinite unitary group; the
relation of the two is \emph{unresolved}: the conversion of Davidson's kernel scalar into the multiplier
\m{\Delta_{\varepsilon}} of the formal text has not been verified.
\end{enumerate}

\subhead{How this note was written.} From the dossier of the unit: one structured proof of ten steps. The
program check of the unit's dossier covered \UnitSteps{} leaf steps of \UnitProofs{} proof, found the statement
at its lines and compared the quotations of the imported declarations with the source. I read the whole Lean file
of the unit, statement and proof, and compared the statement with Theorem~1 argument by argument, in Lean's order:
the index set with its two instance arguments, the signs, the polynomial, the vector, the translation, the group
element, the equation \m{v=\pv{\varepsilon}{p}}, the three natural numbers, and the four hypotheses. In the
background library I read the definitions quoted in Definitions \ref{d:fock} to~\ref{d:weil} (the bicharacter, the
Gaussian, the normalised measure, the space with its reproducing identity, the Heisenberg operators, the inner
product, the polynomial vectors, the multiplier, the operators \m{\weil{\varepsilon}{u}}) and the statements of
the lemmas behind (Q2), (Q4), (Q5) and (Q6); I did not read the proof that polynomial vectors lie in the space
(Q1), nor the proof of the cocycle identity (Q6). The proof above follows the formal one step by step; its
arrangement into three lemmas, the remark on the best~$\nu$ and §\ref{s:use} are mine. For §\ref{s:use} I read the
statement of the consuming lemma and the lines where Theorem~1 is applied.

\head{Appendix A. Concordance with the Lean artifact}
{\fontsize{8.6}{11.6}\selectfont\renewcommand{\leanfont}{\ttfamily\fontsize{7.7}{10}\selectfont\color{leancol}}
\setlength{\LTleft}{0pt}\setlength{\LTright}{\fill}\setlength{\tabcolsep}{0pt}
\begin{longtable}{@{}p{30mm}p{56mm}p{43mm}r@{}}
\textbf{In this note} & \textbf{Lean declaration} & \textbf{File} & \textbf{Lines}\\\midrule\endhead
\ConcordanceRows
\end{longtable}}

\subhead{The statement as it is written in the artifact.}
\quoted{normrealInnerrealHeisOple}
In this text {\leanfont ι} is the index set~$I$, {\leanfont σ} the group element~$u$, {\leanfont P} the
polynomial~$p$, {\leanfont d} the degree~$\dg$, {\leanfont n} the exponent~$\nu$; {\leanfont USign ε} is
$U(\varepsilon)$, {\leanfont polyVector ε P} is \m{\pv{\varepsilon}{p}}, {\leanfont realInner} is the inner
product, {\leanfont realHeisOp y} is \m{\heis{\varepsilon}{y}} and {\leanfont realWeil σ} is
\m{\weil{\varepsilon}{u}}; {\leanfont ‖toLp 2 y‖} is the Euclidean norm~$|y|$, and {\leanfont ‖σ.1.1‖} the operator
norm of the matrix of~$u$. The letters of the formal text were changed where the write-up reserves them for other
objects (the letters \emph{d}, \emph{P} and \emph{n}).

\head{Notation}
\begingroup\fontsize{8.6}{10.9}\selectfont\raggedright
\setlength{\columnsep}{9mm}\setlength{\parskip}{0pt}
\begin{multicols}{2}
\nt{$I$, $|I|$, $i$}{a non-empty finite set, its number of elements, an element of it}
\nt{$\varepsilon$}{a family of signs \m{\varepsilon_i\in\{1,-1\}}, \m{i\in I}}
\nt{$\CC^I$, $|z|$}{the vectors indexed by~$I$; the Euclidean norm}
\nt{$z$, $y$}{vectors of $\CC^I$; $y$ is the translation of Theorem~1}
\nt{$w$}{in this note also a vector of $\CC^I$ (not a finite place, Convention~0.4)}
\nt{$\mathbf{w}$}{an element of the Fock space}
\nt{$w^{*}\varepsilon z$}{the hermitian form \m{\sum_i\overline{w_i}\varepsilon_iz_i}}
\nt{$\dd z$}{Lebesgue measure on $\CC^I$ divided by \m{\pi^{|I|}}}
\nt{$\beta_{\varepsilon}$, $\bich{\varepsilon}{z}{w}$}{the bicharacter \m{\exp(\mathrm{i}\operatorname{Im}(w^{*}\varepsilon z))}}
\nt{$\vac$}{the Gaussian \m{\exp(-|z|^2/2)}, the vacuum vector}
\nt{$\Fock{\varepsilon}$, $\lone{v}$}{the Fock space of signature~$\varepsilon$; its norm}
\nt{$\inn{\varepsilon}{v}{\mathbf{w}}$}{the inner product \m{\int\overline{v}\,\mathbf{w}\,\dd z} of $\Fock{\varepsilon}$}
\nt{$\heis{\varepsilon}{y}$}{the Heisenberg operator of \m{y\in\CC^I}}
\nt{$U(\varepsilon)$, $u$, $u^{*}$}{the unitary group of \m{\operatorname{diag}(\varepsilon)}; an element; its
conjugate transpose}
\nt{$u_{--}$, $\Delta_{\varepsilon}$, $\mult{\varepsilon}{u}$}{the block of~$u$ at the negative signs; the multiplier
\m{\overline{\det u_{--}}}}
\nt{$\opn{u}$}{the operator norm of a matrix for the Euclidean norm}
\nt{$\weil{\varepsilon}{u}$}{the oscillator representation at~$u$}
\nt{$X_i$, $p$}{the variables; a complex polynomial in them}
\nt{$\hol{\varepsilon}{z}$}{$z$ with its coordinates at positive signs conjugated}
\nt{$\pv{\varepsilon}{p}$, $v$}{the Fock vector of~$p$; an element of the Fock space}
\nt{$\dg$}{the degree of the homogeneous polynomial~$p$ (a natural number)}
\nt{$\delta$}{the bound of Theorem~C in the overview; not used for an object of this note}
\nt{$A$}{a natural number with \m{\dg+|I|\le A}}
\nt{$\nu$}{in this note also a natural number: the decay is of order \m{|y|^{-2\nu}} (not a character of the ideles, Convention~0.6)}
\nt{$c$}{a real number; in the proof \m{c=\sqrt{\nu/A+1}}}
\nt{$B$}{a real number; in the proof \m{(|z|^2+|uz+y|^2)/(2(\nu+A))}}
\nt{$T$}{a bounded linear map on the Fock space}
\nt{$j$}{the size of a minor}
\nt{$\alpha$, $\gamma$}{real numbers in elementary inequalities}
\nt{$\ee$, $\pi$}{the base of the natural logarithm; the area of the unit disc}
\end{multicols}
\endgroup

\end{document}
